\documentclass[11pt,a4paper]{article}
\usepackage[margin=1in]{geometry}
\usepackage{amsmath,amssymb}
\usepackage{graphicx}
\usepackage{float}
\usepackage{enumitem}
\usepackage{microtype}
\usepackage{caption}
\usepackage{tabularx}
\usepackage{hyperref}

\captionsetup{font=small,labelfont=bf,skip=6pt}

\setlist[itemize]{leftmargin=1.2cm}
\setlist[enumerate]{leftmargin=1.2cm}

\newcommand{\problembox}[3]{%
  \par\addvspace{0.42em}%
  \noindent
  \setlength{\fboxsep}{7.5pt}%
  \setlength{\fboxrule}{0.95pt}%
  \fbox{%
    \begin{minipage}{\dimexpr\linewidth-2\fboxsep-2\fboxrule\relax}
    \begin{tabularx}{\linewidth}{@{}>{\raggedright\arraybackslash\bfseries}p{2.1em} @{\hspace{0.9em}} >{\raggedright\arraybackslash}X @{\hspace{0.9em}} >{\raggedleft\arraybackslash\bfseries}p{2.4em}@{}}
    #1 & #2 & #3
    \end{tabularx}
    \end{minipage}%
  }%
  \par\addvspace{0.32em}%
}

\title{\textbf{Dynamics of quasiparticles}}
\author{\normalsize Y23 (2023) --- English translation}
\date{}
\begin{document}
\maketitle

Solids (such as metals) are a crystal lattice of positively charged ions between which electrons move. Electrons interact with the crystal lattice and with each other, so to accurately describe the properties of a solid it is necessary to solve an extremely complex multiparticle problem. However, in some cases it turns out that the entire influence of interaction comes down to the fact that instead of electrons in the crystal, particles weakly interacting with each other propagate, the equations of motion of which differ significantly from the equations for free electrons. Such particles that practically do not interact with each other are called quasiparticles. For simplicity, they are also usually called electrons. The properties of quasiparticles are determined by the parameters of the crystal lattice and the magnitude of the interaction between electrons and ions. For an electron (quasiparticle), the dependence of energy on momentum $\varepsilon(\vec{p})$ can have a very complex form. Then the connection between velocity and momentum is given by the relations:
$$ v_x = \frac{\partial \varepsilon}{\partial p_x}, \quad v_y = \frac{\partial \varepsilon}{\partial p_y}, \quad v_z = \frac{\partial \varepsilon}{\partial p_z}. $$
Speed is understood in the usual sense as the derivative of the corresponding coordinate with respect to time. The expression for the force acting on an electron in external electric and magnetic fields is the same as in the case of motion in a vacuum, and the equations of motion have the form:
$$ \frac{d\vec{p}}{dt} = \vec{F}. $$
The electron charge is $-q$ with $q > 0$. In Part C, you can study another type of quasiparticle — phonons, which are vibrations of the crystal lattice.

\section*{Part A. Bloch oscillations (2 points)}

\problembox{A1}{Let the dependence of electron energy on momentum have the form $$\varepsilon(p) = \varepsilon_0 (3 - \cos(ap_x) - \cos(ap_y) - \cos(ap_z)).$$ Find its acceleration at the moment of time when the momentum is $p_x = \frac{\pi}{2 a}$, $p_y = \frac{\pi}{3 a}$, $p_z = 0$, and the components of the force acting on the electron are $F_x$, $F_y$, $F_z$.}{0.70}

\problembox{A2}{Let an electron with $\varepsilon(\vec{p})$ from A1 move in a constant electric field of magnitude $E$ directed along the $z$ axis. The initial coordinates and projections of the momentum are equal to zero. Find its law of motion $z(t)$.}{0.80}

\problembox{A3}{The concentration of free electrons in a certain substance is $n$, and the dependence $\varepsilon(p)$ is as in A1. At the initial moment of time, a constant electric field $E$ is turned on, directed along the $z$ axis. Find the dependence of the current density $j$ on time.}{0.50}

\section*{Part B. Motion in a magnetic field (4 points)}

In this part, we will consider the motion of an electron in a constant magnetic field $\vec{B}$.

\problembox{B1}{Let the dependence of energy on momentum $\varepsilon(\vec{p})$ be arbitrary. Prove that the energy $\varepsilon$ and the projection of momentum on the direction of the magnetic field are conserved.}{0.50}

We will further assume that the magnetic field of magnitude $B$ is directed along the $z$ axis, and the electron moves in the $xy$ plane.

\problembox{B2}{Let the dependence of energy on momentum have the form $$\varepsilon = \frac{p_x^2}{2 m_x} + \frac{p_y^2}{2 m_y} .$$ The electron energy is $\varepsilon$. Find the frequency of motion $\omega$. Depict the trajectory of the particle on the $xy$ plane, indicate the characteristic dimensions and direction of motion. Place the origin of coordinates at the center of symmetry of the trajectory.}{1.10}

Let now, in addition to the magnetic field, the electron from B2 be subject to an electric field $E$ directed along the $x$ axis. The initial values of the coordinate and momentum of the electron are zero.

\problembox{B3}{Find the hodograph of the electron's momentum vector (the trajectory along which the end of the momentum vector moves if it is plotted from a single point). Draw it on the $p_x, p_y$ plane. Please indicate characteristic dimensions.}{0.60}

\problembox{B4}{Find the law of motion $x(t)$, $y(t)$.}{0.40}

The electric field is again zero. The dependence of electron energy on momentum has the form:
$$\varepsilon = \varepsilon_0(1 - \cos(ap_x)) + \frac{p_y^2}{2 m_y}.$$

\problembox{B5}{If the energy is greater than the minimum value, the trajectories on the plane become open, and the particle can travel indefinitely in some direction. Find this minimum energy $\varepsilon_{min}$. At energy $\varepsilon > \varepsilon_{min}$, draw the trajectory and indicate the characteristic geometric dimensions.}{1.40}

\section*{Part C. Phonon decay (4 points)}

Another type of quasiparticle found in solid state physics is phonons, which represent vibrations of a crystal lattice. Their properties are similar to photons, but due to lattice characteristics, the dependence of energy on momentum is nonlinear:
$$\varepsilon(p) = up + \alpha p^3.$$
Here $u$ is the speed of sound in the crystal, and the constant $\alpha$ sets the magnitude of dispersion. In this part, we assume the second term is much smaller than the first: $\alpha p^3 \ll up$. Unlike a photon, a phonon under certain conditions can decay into two phonons.

Let the momentum of the original phonon be $p$, and it decays into two phonons with momenta $\vec{q}_1$ and $\vec{q}_2$ which move at angles $\theta_1, \theta_2$ relative to the direction of motion of the original phonon.

\begin{figure}[H]
    \centering
    \includegraphics[width=0.4\textwidth]{T2_EN_assets/fig1.jpg}
    \caption{Figure 1.}
\end{figure}

\problembox{C1}{First, neglect dispersion and consider the case $\alpha = 0$. At what angles $\theta_1, \theta_2$ is decay possible?}{0.50}

Next we consider the case $\alpha \ne 0$. It turns out that for a small momentum $p$ of the initial phonon, the angles $\theta_1, \theta_2$ are also small.

\problembox{C2}{Express the modulus of the second phonon momentum $q_2$ in terms of $p, q_1, \theta_1$. Obtain the expression in the limit of small $\theta_1$ up to second order in $\theta_1$.}{0.50}

\problembox{C3}{The dependence of the emitted phonon momentum on the angle at small angles has the form: $$q_1(\theta_1) = A + B \theta_1^2.$$ Determine the constants $A, B$. Express them in terms of $p, u, \alpha$.}{1.20}

\problembox{C4}{At what sign of the constant $\alpha$ is decay possible?}{0.30}

\problembox{C5}{Find the maximum angle $\theta_{max}$ at which a phonon can be emitted. Under what condition on $p$ can the phonon momentum be considered small?}{0.50}

\problembox{C6}{Express the angle $\theta_2$, under which the second phonon is emitted, in terms of $\theta_1, p, u, \alpha$. Consider the angles small and account for first-order contributions in $\theta_1, \theta_2$.}{1.00}

\vspace{1em}
\noindent\textit{Source:} \url{https://pho.rs/p/3148}

\end{document}